\documentclass[12pt, a4paper]{article}

\usepackage{amsmath, amssymb}
\usepackage{fontspec}
\setmainfont{FreeSerif}
\usepackage{geometry}
\geometry{margin=1in}
\usepackage[colorlinks=true, linkcolor=blue, urlcolor=blue, citecolor=blue]{hyperref}

\setcounter{secnumdepth}{0} 

\begin{document}

\tableofcontents
\vspace{2em}
\hrule
\vspace{2em}

\section{Հիմնական սահմանումներ և բաց բազմություններ}

$\mathbb{R}^m$ տարածության մեջ $x \in \mathbb{R}^m$ կետի $\varepsilon$-շրջակայք (կամ \textbf{բաց գունդ}) կոչվում է հետևյալ բազմությունը.

$$B(x, \varepsilon) = \{y \in \mathbb{R}^m : \|x - y\| < \varepsilon\}$$

\vspace{1em}
\hrule
\vspace{1em}

Դիցուք ունենք $F \subseteq \mathbb{R}^m$ բազմությունը։

$a \in F$ կետը կոչվում է $F$ բազմության ներքին կետ, եթե գոյություն ունի այնպիսի $\varepsilon > 0$ թիվ, որ՝

$$B(a, \varepsilon) \subseteq F$$

\vspace{1em}
\hrule
\vspace{1em}

\textbf{Բաց բազմություն:} $F \subseteq \mathbb{R}^m$ բազմությունը կոչվում է բաց, եթե նրա յուրաքանչյուր կետ հանդիսանում է ներքին կետ, այսինքն՝

$$\forall x \in F, \exists \varepsilon > 0 \text{ s.t. } B(x, \varepsilon) \subseteq F$$

\textbf{Փակ բազմություն}: $F$ բազմությունը կոչվում է փակ, եթե $F'\subset F$ ($F'$ - կուտակման կետերի բազմություն)

\vspace{1em}
\hrule
\vspace{1em}

\subsection{Թեորեմ: Բաց և փակ բազմությունների կապը}

Որպեսզի $F \subseteq \mathbb{R}$ բազմությունը լինի փակ, անհրաժեշտ է և բավարար, որ նրա լրացումը՝ $V = \mathbb{R} \setminus F$ բազմությունը, լինի բաց։

\textbf{Ապացույց}

\begin{enumerate}
    \item Անհրաժեշտություն ($\implies$)
\end{enumerate}

\textbf{Պայման:} Ենթադրենք $F$-ը փակ բազմություն է։

\textbf{Պահանջ}: Ապացուցել, որ $V = \mathbb{R} \setminus F$-ը բաց է։

Վերցնենք կամայական $x \in V$ կետ։ Քանի որ $V = \mathbb{R} \setminus F$, ապա $x \notin F$։

Քանի որ $F$-ը փակ բազմություն է, այն պարունակում է իր բոլոր սահմանային (կուտակման) կետերը ($F' \subseteq F$)։

Քանի որ $x \notin F$, ապա $x$-ը չի կարող լինել $F$-ի կուտակման կետ ($x \notin F'$)։

Ըստ կուտակման կետի սահմանման՝ դա նշանակում է, որ գոյություն ունի $x$ կետի այնպիսի $\varepsilon$-շրջակայք, որը չի պարունակում $F$-ի տարրեր (բացի, հնարավոր է, հենց $x$-ից, բայց $x \notin F$)։

Այսինքն՝ $\exists \varepsilon > 0$, որ $(x - \varepsilon, x + \varepsilon) \cap F = \emptyset$։

Սա նշանակում է, որ ամբողջ շրջակայքն ընկած է լրացման մեջ. $(x - \varepsilon, x + \varepsilon) \subseteq V$։

Հետևաբար, $x$-ը $V$ բազմության ներքին կետ է։ Քանի որ $x$-ը կամայական էր, $V$-ն բաց բազմություն է։

\begin{enumerate}
    \setcounter{enumi}{1}
    \item Բավարարություն ($\impliedby$)
\end{enumerate}

\textbf{Պայման:} Ենթադրենք $V = \mathbb{R} \setminus F$-ը բաց բազմություն է։

\textbf{Պահանջ:} Ապացուցել, որ $F$-ը փակ է (այսինքն՝ $\forall x \in F' \implies x \in F$)։

Ապացուցենք հակասող ենթադրությամբ։ Ենթադրենք գոյություն ունի $x \in F'$ կուտակման կետ, որը չի պատկանում $F$-ին ($x \notin F$)։

Այդ դեպքում $x \in V$։

Քանի որ $V$-ն բաց է, ապա ըստ սահմանման $x$-ը ներքին կետ է, այսինքն՝ $\exists \varepsilon > 0$, որ $(x - \varepsilon, x + \varepsilon) \subseteq V$։

Սա նշանակում է, որ այդ շրջակայքը $F$ բազմության հետ ընդհանուր տարր չունի. $(x - \varepsilon, x + \varepsilon) \cap F = \emptyset$։

Սակայն սա հակասում է այն փաստին, որ $x$-ը $F$ բազմության կուտակման կետ է (կուտակման կետի ցանկացած շրջակայք պետք է պարունակի $F$-ի տարրեր)։

Հետևաբար, մեր ենթադրությունը սխալ էր, և $x$-ը պետք է պատկանի $F$-ին։ Ուստի $F$-ը փակ բազմություն է։

$\blacksquare$

\vspace{1em}
\hrule
\vspace{1em}

\textbf{Եզրային կետի սահմանումը}: $a \in \mathbb{R}^m$ կետը կոչվում է $F \subseteq \mathbb{R}^m$ բազմության եզրային կետ, եթե $a$ կետի կամայական շրջակայքում կան թե՛ $F$ բազմությանը պատկանող և թե՛ $F$-ին չպատկանող կետեր:

$F$ բազմության բոլոր եզրային կետերի բազմությունը նշանակվում է $\partial F$ սիմվոլով:

Եթե $F$ բազմությունը բաց տիրույթ է, ապա $F \cup \partial F$ բազմությունը փակ տիրույթ է:

\vspace{1em}
\hrule
\vspace{1em}

\textbf{Սահմանում։}
$\mathbb{R}^m$ տարածության $a$ և $b$ կետերը միացնող հատված ասելով կհասկանանք [0, 1]֊ի պատկերը

$$\varphi(t)= (1-t)a + tb = a + (b-a)t, \ t \in [0, 1]$$
արտապատկերման ժամանակ։

$a$֊ն կոչվում է սկբնակետ, $b$֊ն վերջնակետ։ $a, b \in \mathbb{R}^m$

\subsection{Ուռուցիկ բազմություններ}

\textbf{Սահմանում։} $L \in \mathbb{R}^m$ բազմությունը կոչվում է բեկյալ, եթե գոյություն ունի կետերի այնպիսի $\{x_0, x_1, \dots, x_n\} \subseteq \mathbb{R}^m$ բազմություն, որ $L$ բազմությունը հանդիսանում է այդ կետերը միացնող հատվածների միավորումը:

$$L = \bigcup_{i=0}^{n-1} [x_i, x_{i+1}]$$

Այսինքն՝ բեկյալը բաղկացած է վերջավոր թվով իրար հաջորդող հատվածներից, որտեղ յուրաքանչյուր հատվածի վերջնակետը հաջորդ հատվածի սկզբնակետն է:

\textbf{Սահմանում։} $L \in \mathbb{R}^m$ բազմությունը կոչվում է կապակցված, եթե $\forall a, b \in F \quad \exists L \subset F, \text{  s.t.  }  x_0 = a, x_n = b$

$F \subseteq \mathbb{R}^m$ բազմությունը կոչվում է ուռուցիկ, եթե նրան պատկանող ցանկացած երկու կետերի համար նրանց միացնող հատվածը նույնպես պատկանում է այդ բազմությանը.

$$\forall a, b \in F \implies [a, b] \subseteq F$$

\subsection{Կոմպակտ բազմություններ և Բորելի լեմման}

\textbf{Կոմպակտ բազմության սահմանումը:}
$F \subseteq \mathbb{R}^m$ բազմությունը կոչվում է կոմպակտ, եթե $F$-ը ծածկող կամայական բաց բազմությունների ընտանիքից կարելի է առանձնացնել վերջավոր ենթածածկույթ։

Ավելի պարզ ասած՝ եթե դու ունես անվերջ թվով բաց բազմություններ, որոնց միավորումը պարունակում է $F$-ը, ապա կոմպակտության շնորհիվ դու կարող ես ընտրել դրանցից ընդամենը մի քանիսը (վերջավոր քանակ), որոնք դեռևս ամբողջությամբ կծածկեն քո բազմությունը։

Եթե տրված է կամայական բաց ծածկույթ $\{G_\alpha \mid \alpha \in I\}$ այնպիսին, որ $F \subseteq \bigcup_{\alpha \in I} G_\alpha$, ապա գոյություն ունեն վերջավոր թվով ինդեքսներ $\alpha_1, \alpha_2, \dots, \alpha_n$ այնպես, որ՝

$$F \subseteq \bigcup_{i=1}^{n} G_{\alpha_i}$$

\vspace{1em}
\hrule
\vspace{1em}

\subsection{1. Կոմպակտ բազմության սահմանափակությունը}

\textbf{Թեորեմ:} Եթե $F \subset \mathbb{R}^m$ բազմությունը կոմպակտ է, ապա այն սահմանափակ է։

\textbf{Ապացույց:}

Դիցուք $F \subset \mathbb{R}^m$ և $F$-ը կոմպակտ բազմություն է։

Դիտարկենք կոորդինատների սկզբնակետում կենտրոն ունեցող և $n$ բնական շառավղով բաց գնդերի հաջորդականությունը՝ $B(0, n)_{n \in \mathbb{N}}$։

Քանի որ շառավիղն անվերջ աճում է, ապա այդ գնդերի միավորումը ծածկում է ամբողջ $\mathbb{R}^m$ տարածությունը՝

$$\mathbb{R}^m = \bigcup_{n \in \mathbb{N}} B(0, n)$$

Քանի որ $F \subset \mathbb{R}^m$, ապա $F$ բազմությունը նույնպես ծածկվում է այդ գնդերի միջոցով՝

$$F \subset \bigcup_{n \in \mathbb{N}} B(0, n)$$

Սա նշանակում է, որ $\{B(0, n) \mid n \in \mathbb{N}\}$ ընտանիքը $F$ բազմության բաց ծածկույթ է։

Օգտվելով $F$-ի կոմպակտության սահմանումից՝ այս անվերջ ծածկույթից կարելի է առանձնացնել վերջավոր ենթածածկույթ։

Այսինքն՝ գոյություն ունեն վերջավոր թվով բնական թվեր $n_1, n_2, \dots, n_k$ այնպիսին, որ՝

$$F \subset \bigcup_{i=1}^{k} B(0, n_i)$$

Քանի որ գնդերը նույն կենտրոնն ունեն, ապա դրանց միավորումը հավասար է ամենամեծ շառավղով գնդին՝

$$\bigcup_{i=1}^{k} B(0, n_i) = B(0, \max_{1 \le i \le k} \{n_i\})$$

Հետևաբար, $F$ բազմությունն ամբողջությամբ ընկած է մի վերջավոր շառավղով գնդի մեջ, ինչը նշանակում է, որ $F$-ը սահմանափակ է։

\subsection{2. Կոմպակտ բազմության փակությունը}

\textbf{Թեորեմ:} Եթե $F \subset \mathbb{R}^m$ բազմությունը կոմպակտ է, ապա այն փակ է։

\textbf{Ապացույց} (հակասող ենթադրությամբ):

Դիցուք $x_0$-ն $F$ բազմության կուտակման կետ է։

Ենթադրենք հակառակը՝ $x_0 \notin F$։

Դիտարկենք հետևյալ բաց բազմությունների ընտանիքը՝

$$F_n = \mathbb{R}^m \setminus \overline{B(x_0, 1/n)}, \quad n \in \mathbb{N}$$

Այս ընտանիքի միավորումը ծածկում ամբողջ տարածությունը՝ բացի $x_0$ կետից՝

$$\bigcup_{n \in \mathbb{N}} F_n = \mathbb{R}^m \setminus \{x_0\}$$

Քանի որ $x_0 \notin F$, ապա $F \subset \bigcup_{n \in \mathbb{N}} F_n$, այսինքն՝ $\{F_n\}$ ընտանիքը $F$ բազմության բաց ծածկույթ է։

Օգտվելով $F$-ի կոմպակտությունից՝ առանձնացնենք վերջավոր ենթածածկույթ $n_1, n_2, \dots, n_k$ ինդեքսներով՝

$$F \subseteq \bigcup_{i=1}^{k} F_i$$

Քանի որ $F_n$ բազմությունները ներդրված են միմյանց մեջ, ապա դրանց միավորումը հավասար է $F_{n_{\max}}$-ին՝

$$F \subseteq \mathbb{R}^m \setminus \overline{B(x_0, 1/n_{\max})}$$

Սա նշանակում է, որ $x_0$ կետի $1/n_{\max}$ շառավղով շրջակայքը չի պարունակում $F$ բազմության ոչ մի տարր՝

$$B(x_0, 1/n_{\max}) \cap F = \emptyset$$

Սա հակասում է այն պայմանին, որ $x_0$-ն $F$-ի կուտակման կետ է։

Հետևաբար, $x_0 \in F$, ինչը նշանակում է, որ $F$ բազմությունը փակ է։

\section{Բորելի լեմմա}

$\mathbb{R}^m$ տարածության մեջ $F$ բազմությունը կոմպակտ է այն և միայն այն դեպքում, երբ այն փակ է և սահմանափակ:

\textbf{Ապացույց} ($m=2$ դեպքի համար)

Նախորդ հատվածներում արդեն ապացուցվել է անհրաժեշտությունը (եթե կոմպակտ է, ապա փակ է և սահմանափակ): Այժմ ապացուցենք բավարարությունը:

\textbf{Պայման:} Դիցուք $F \subseteq \mathbb{R}^2$ բազմությունը փակ է և սահմանափակ:

\textbf{Պահանջ:} Ապացուցել, որ $F$-ը կոմպակտ է, այսինքն՝ նրա կամայական $\{G_\alpha\}_{\alpha \in I}$ բաց ծածկույթից կարելի է առանձնացնել վերջավոր ենթածածկույթ:

\textbf{Ենթադրություն (հակասող):} Ենթադրենք $F$-ը չի կարելի ծածկել $\{G_\alpha\}$ համակարգի վերջավոր թվով տարրերով:

Քառակուսիների կառուցումը: Քանի որ $F$-ը սահմանափակ է, այն զետեղված է ինչ-որ $A_1 = [a_1, b_1] \times [c_1, d_1]$ ուղղանկյան (կամ քառակուսու) մեջ:

Տրոհենք $A_1$ ուղղանկյունը 4 հավասար մասերի:

Ըստ մեր ենթադրության՝ այդ մասերից առնվազն մեկը (նշանակենք $A_2$) այնպիսին է, որ $A_2 \cap F$ բազմությունը նույնպես հնարավոր չէ ծածկել վերջավոր ենթածածկույթով:

Շարունակելով այս պրոցեսը՝ կստանանք ներդրված ուղղանկյունների հաջորդականություն՝ $A_1 \supset A_2 \supset \dots \supset A_n \dots$, որոնց կողմերը ձգտում են զրոյի ($|a_n - b_n| = \frac{b-a}{2^{n-1}}$ և $|c_n - d_n| = \frac{d-c}{2^{n-1}}$), և $A_n \cap F$ բազմություններից ոչ մեկը չի ծածկվում վերջավոր ենթածածկույթով:

\textbf{Ընդհանուր կետի գոյությունը:}

Ըստ Կանտորի թեորեմի (ներդրված հատվածների սկզբունքի)՝ գոյություն ունի միակ $(a, c)$ կետ, որը պատկանում է բոլոր $A_n$ ուղղանկյուններին:

Քանի որ $F$-ը փակ է և $A_n \cap F \neq \emptyset$, ապա այդ սահմանային կետը պատկանում է $F$-ին՝ $(a, c) \in F$։

\textbf{Հակասության ստացումը:}

Քանի որ $\{G_\alpha\}$-ն ծածկում է $F$-ը, ապա գոյություն ունի ինչ-որ $\alpha_0$, որ $(a, c) \in G_{\alpha_0}$:

Քանի որ $G_{\alpha_0}$-ն բաց բազմություն է, ապա $(a, c)$ կետը դրա մեջ է մտնում իր ինչ-որ շրջակայքով:

Բավականաչափ մեծ $n$ ինդեքսի համար $A_n$ ուղղանկյունն ամբողջությամբ կհայտնվի այդ շրջակայքում, հետևաբար՝ $A_n \subseteq G_{\alpha_0}$:

Սա նշանակում է, որ $A_n \cap F$ բազմությունը ծածկվում է ընդամենը մեկ բաց բազմությամբ ($G_{\alpha_0}$), ինչը հակասում է մեր ենթադրությանը, թե այն հնարավոր չէ ծածկել վերջավոր թվով բազմություններով:

Եզրակացություն: Ենթադրությունը սխալ էր, հետևաբար $F$-ը կոմպակտ է:

$\blacksquare$

\end{document}